\sscp{Preposed possessor}{13-03-04}
In \srf{01-05}
and \srf{02-03-01-01} we discuss how the order of the possessor
in possessive constructions having \it{preposed possessors} has
been used by a number of scholars in the past to distinguish
Austronesian languages in the west that have post-posed possessors (e.g.
\ili{Malay} \it{kepala \tbf{saya}} ‘my head’) from those in the \ili{east}
that have preposed possessors (e.g.
\ili{Kotos} \ili{Amarasi} \it{\tbf{au} ʔnakak} ‘my head’).
We noted that having preposed possessors is typologically unusual
for SVO languages that have prepositions.
Subsequently the presence of preposed possessors in so many Austronesian
languages in the target region has been shown to be due to intense
contact over a long period of time with Papuan languages which
are typologically SOV and have preposed possessors commonly associated with that typology.
\ili{As} a typological areal feature resulting from contact,
the presence of preposed possessors cannot be used as a basis for linguistic subgrouping.

The distribution of the preposed possessor only partially mirrors
the Blust line (see \frf{02-16}).
Languages in the western,
central and southern parts of the \tsc{\ili{Bima}-Lembata} subgroup
do not have the preposed possessor,
but their historical phonologies set them apart from subgroups
further west, and they share some lexical features
and other things related to contact that link them loosely with
languages to the \ili{east} that do (see \crf{ch:BL}).
\ili{Banggai} and \tsc{Taliabu} have preposed possessors,
but their historical phonologies group them westward with the
\tsc{Celebic} languages of Sulawesi (see \crf{ch:Tal}).